\documentclass[11pt]{article}
    \usepackage{amsmath}
    \usepackage{amssymb}
    \title{\textbf{Forelle} 1.0}
    \author{Brass Glass}
    \date{}
    
    \addtolength{\topmargin}{-3cm}
    \addtolength{\textheight}{3cm}
\begin{document}
\maketitle
\thispagestyle{empty}

\section{Overview}
Welcome to version 1.0 of \textbf{Forelle} — a GtkSourceView\footnote{https://gitlab.gnome.org/GNOME/gtksourceview} (GNOME-based) styling sheet! This style focusses on maintaining a warm palette, using a cream background, teal sidebar, and pink \& green texts.\\
This syling sheet is inspired off of SurrealEmber’s “Poisson 23” from the platform Lospec\footnote{https://lospec.com/palette-list/poisson-23}. Credit to him for this amazing palette!

\section{Showcase}
Following are some showcases of Forelle using \LaTeX\:features.\\
Do this verbatim!
\begin{verbatim}
This text is done verbatim.
If you look to the side, these texts display as typewriter!
Anything is done verbatim — for example, \LaTeX is not formatted.
\end{verbatim}
A multi-line proof detailing the bijection which exists between the Naturals and the Integers via the Cantor-Bernstein-Schröeder Theorem.
\begin{align*}
f \colon \mathbb{N} &\to \mathbb{Z}   &  g \colon \mathbb{Z} &\to \mathbb{N}\\
n &\mapsto n  &  
n &\mapsto \begin{cases}
	3^n & n \equiv 0 \pmod{2}\\
	2^{-n} & n \equiv 1 \pmod{2}
\end{cases}\\
\end{align*}
Therefore, there must exist a bijection between the naturals and integers.\\\\
Lastly, we will test a simple in-line math. If $25 \cdot 34 \equiv 83 \pmod{16}$, then Gauss didn’t implement the modulo into mathematics properly.\\
% Oh, and this is a comment! It does not show in the text.
% This screenshot was taken using Gummi, an offline LaTeX editor. Credit to its creätors!
% The font in this screenshot is monofur, a rounded monospaced geometric font designed by Tobias Benjamin Köhler. Credit to him!

\section{Closing}
This is all of Forelle for now! Any feedback is appreciated, and focussed compatability will expand for more GtkSourceView-supportive applications.\\
Enjoy!

\end{document}

